\medskip

Le but de cet exercice est d'étudier les convergences de deux suites vers une même limite.

\medskip

\textbf{Partie A}

\medskip

On considère la fonction $f$ définie sur $[2~;~+\infty[$ par 

\[f(x) = \sqrt{3x - 2}.\]

\medskip

\begin{enumerate}
\item Justifier les éléments du tableau de variations ci-dessous :











\[\begin{tablvar}[stretch=1.2,intervalwidth=8em]{1}
\hline
x & 2 && +\infty\\
\hline
\variations{\mil{f(x)} & \bas{2} && \haut{+\infty}}
\hline
\end{tablvar}\]

On admet que la suite $\left(u_n\right)$ vérifiant $u_0 = 6$ et, pour tout $n$ entier naturel, $u_{n+1} = f\left(u_n\right)$ est bien définie.
\item 
	\begin{enumerate}
		\item Démontrer par récurrence que, pour tout $n$ entier naturel : $2 \leqslant u_{n+1} \leqslant u_n \leqslant 6$.
		\item En déduire que la suite $\left(u_n\right)$ converge.
	\end{enumerate}
\item On appelle $\ell$ la limite de $\left(u_n\right)$.

On admet qu'elle est solution de l'équation $f(x) = x$. Déterminer la valeur de $\ell$.

\item On considère la fonction \texttt{rang} écrite ci-dessous en langage Python.

On rappelle que sqrt($x$) renvoie la racine carrée du nombre $x$.

\begin{center}
\begin{ttfamily}
\begin{tabular}{>{\cellcolor{lightgray!50}}c l}
1&from math import *\\
2&\\
3&def rang(a) :\\
4&\quad  u = 6\\
5&\quad n=0\\
6& \quad while u > = a :\\
7&\qquad u = sqrt(3*u - 2)\\
8&\qquad n = n+1\\
9& \quad return n\\
\end{tabular}
\end{ttfamily}
\end{center}

	\begin{enumerate}
		\item Pourquoi peut-on affirmer que \texttt{rang}(2.000001) renvoie une valeur ?
		\item Pour quelles valeurs du paramètre $a$ l'instruction \texttt{rang}($a$) renvoie-t-elle un résultat ?
	\end{enumerate}
\end{enumerate}

\bigskip

\textbf{Partie B}

\medskip

On admet que la suite $\left(v_n\right)$ vérifiant $v_0 = 6$ et, pour tout $n$, entier naturel, $v_{n+1} = 3 - \dfrac{2}{v_n}$ est bien définie.

\medskip

\begin{enumerate}
\item Calculer $v_1$.
\item Pour tout $n$ entier naturel, on admet que $v_n \ne 2$ et on pose :

\[w_n = \dfrac{v_n - 1}{v_n - 2}\]

	\begin{enumerate}
		\item Démontrer que la suite $\left(w_n\right)$ est géométrique de raison 2 et préciser son premier terme $w_0$.
		\item On admet que, pour tout $n$ entier naturel,

\[w_n - 1 = \dfrac{1}{v_n - 2}.\]

En déduire que, pour tout $n$ entier naturel,

\[v_n = 2 + \dfrac{1}{1,25 \times 2^n - 1}\]

		\item Calculer la limite de $\left(v_n\right)$.
	\end{enumerate}
\item Déterminer le plus petit entier naturel $n$ pour lequel $v_n < 2,01$ en résolvant
l'inéquation.
\end{enumerate}

\bigskip

\textbf{Partie C}

\medskip

À l'aide des parties précédentes, déterminer le plus petit entier $N$ tel que pour tout $n \geqslant N$, les termes $v_n$ et $u_n$ appartiennent à l'intervalle ]1,99~;~2,01[.

\bigskip
